\scanpage{45}
\mspage{7}
% archive sheet number: 44

faut et il suffit qu'on ait aussi les 24 relations (7). Quand \inserted{les} $\deleted{\unsure{\operatorname{Aut}}}\operatorname{End}(V_i)$ sont commutatifs (rappelons qu'ils sont iso entre eux) alors \deleted{les \unsure{isom|isoms}} \inserted{les seules} \deleted{$\operatorname{Aut} V_i$} relations (2), (3) impliquent que l'on a un syst. trans. d'iso. entre les \inserted{anneaux} $\deleted{\unsure{\operatorname{Aut}}}\operatorname{End}(V_i)$, soit $\deleted{\unsure{A}}\unsure{k|h}$ \deleted{\unsure{est}}\unsure{cet} anneau. Alors les relations 5, 6, 7 s'interprètent comme des relations dans \deleted{\unsure{E|Ens}} $k$

\[
\left\{\begin{aligned}
\unsure{\rho|f}_{\deleted{\unsure{ij}}j i k\unsure{l|\ell}}&=\rho_{ijlk}=\rho_{ijkl}^{-1} \qquad (\Rightarrow \rho_{\deleted{\unsure{ij}}jilk}=\rho_{ijkl})\\
\deleted{\unsure{\rho}}\rho_{ikjl}&=1-\rho_{ijkl}
\end{aligned}\right. \tag*{(8)}
\]

En d'autres termes, faisant opérer $\unsure{\mathfrak S|\mathscr S}_3$ sur l'ens \inserted{$U_{\deleted{\unsure{k}}}(k)$} des $\lambda\in k$ tels que $\lambda$ et $1-\lambda$ sont inversibles, par

\[
\begin{aligned}
\sigma_{12}(\lambda)&=\sigma_{34}(\lambda)=\lambda^{-1}\\
\sigma_{23}(\lambda)&=1-\lambda
\end{aligned}
\]

\margin{donc $\mathfrak S_4$ opère par l'intermédiaire du quotient $\mathfrak S_3$ de $\mathfrak S_4$}

\deleted{\unsure{avec} $\unsure{\rho}$} l'application \deleted{$\rho_{ijkl}$}

\[
\underset{\text{torseur sous }\mathfrak S_4}{\operatorname{Rep}(I)}\longrightarrow \deleted{\unsure{k}}U(k)
\qquad (i,j,k,l)\unsure{\text{---}|\longrightarrow}\rho_{ijkl}
\]

est comp. aux $\mathfrak S_4$
